\documentclass[10pt]{extarticle}
\usepackage[a4paper,margin=1.5cm]{geometry}
\usepackage[T1]{fontenc}
\usepackage{amsmath,amssymb,mathtools}
\usepackage{graphicx}
\usepackage[export]{adjustbox}
\usepackage{multicol}
\usepackage{enumitem}
\usepackage{booktabs,array,tabularx}
\setlength{\columnsep}{0.7cm}
\setlength{\parindent}{0pt}
\setlength{\parskip}{3pt}
\pagestyle{empty}
\makeatletter
\newcommand{\qfigdef}[4]{\@namedef{qf@file@#1}{#2}\@namedef{qf@w@#1}{#3}\@namedef{qf@wide@#1}{#4}}
\newcommand{\qf@check}[1]{\@ifundefined{qf@file@#1}{\PackageError{score2}{Unknown figure id #1}{}}{}}
\newcommand{\qf@img}[4][]{\includegraphics[width=#2\dimexpr\@nameuse{qf@w@#4}pt\relax,max width=#3,max height=0.45\textheight,#1]{\@nameuse{qf@file@#4}}}
\newif\ifqf@full
\newif\ifqf@boxed
\newsavebox{\qf@box}
\newcommand{\qf@open}{\global\qf@fulltrue\par\global\qf@boxedtrue
  \setbox\qf@box\vbox\bgroup\hsize\textwidth\linewidth\textwidth}
\newcommand{\qf@emit}[1]{\par\egroup\global\qf@boxedfalse
  \ifdim\dimexpr\ht\qf@box+\dp\qf@box\relax<#1\textheight\noindent\usebox{\qf@box}\par\else\unvbox\qf@box\fi}
\newcommand{\qfullsplit}{\ifqf@boxed\qf@emit{0.85}\fi}
\newcommand{\qf@close}{\ifqf@boxed\qf@emit{0.45}\fi\global\qf@fullfalse}
\newcommand{\qfullbegin}{\end{multicols}\qf@open}
\newcommand{\qfullend}{\qf@close\begin{multicols}{2}}
\newcommand{\qfullstart}{\qf@open}
\newcommand{\qfulldone}{\qf@close}
\newcommand{\qfig}[2][1]{\qf@check{#2}%
  \ifqf@full\qf@blk{#1}{#2}\else
  \ifnum\@nameuse{qf@wide@#2}=1 \expandafter\@firstoftwo\else\expandafter\@secondoftwo\fi
  {\end{multicols}\begin{center}\qf@img{#1}{\textwidth}{#2}\end{center}\begin{multicols}{2}}%
  {\qf@blk{#1}{#2}}\fi}
\newcommand{\qf@blk}[2]{\par\smallskip{\centering\qf@img{#1}{\linewidth}{#2}\par}\smallskip}
\newcommand{\qfigi}[2][1]{\qf@check{#2}\qf@img[valign=c]{#1}{\linewidth}{#2}}
\makeatother
\newcommand{\papersec}[1]{\par\bigskip{\centering\Large\bfseries PAPER-#1\par}\medskip}
\newcommand{\partsec}[2]{\par\medskip\noindent{\large\bfseries PART-#1: #2}\par\smallskip}
\newcommand{\sect}[1]{\par\smallskip\noindent\textbf{#1}\par}
\newcommand{\qn}[1]{\par\smallskip\noindent\textbf{#1.}~}
\newcommand{\paraq}[2]{\par\smallskip\noindent\textbf{Paragraph for Questions #1 and #2}\par}
\newenvironment{qoptions}{\begin{itemize}[leftmargin=1.8em,itemsep=1pt,topsep=2pt,label={}]}{\end{itemize}}
\newcommand{\opt}[2]{\item[(#1)] #2}
\newcolumntype{L}[1]{>{\raggedright\arraybackslash}p{#1}}
\newcommand{\fwstart}{\end{multicols}\par\noindent\begin{minipage}{\textwidth}\small\setlength{\tabcolsep}{4pt}\renewcommand{\arraystretch}{1.5}}
\newcommand{\fwend}{\end{minipage}\par\medskip\begin{multicols}{2}}
% figure registry (generated by score2_tools.py finalize)
\qfigdef{p1_phy_s1_q01_1}{SRG (Enthusiast) - Test 07 (27 Jun 2021)-img/p1_phy_s1_q01_1.png}{137.3}{0}
\qfigdef{p1_phy_s1_q02_1}{SRG (Enthusiast) - Test 07 (27 Jun 2021)-img/p1_phy_s1_q02_1.png}{107.4}{0}
\qfigdef{p1_phy_s1_q03_1}{SRG (Enthusiast) - Test 07 (27 Jun 2021)-img/p1_phy_s1_q03_1.png}{124.5}{0}
\qfigdef{p1_phy_s1_q05_2}{SRG (Enthusiast) - Test 07 (27 Jun 2021)-img/p1_phy_s1_q05_2.png}{116.3}{0}
\qfigdef{p1_phy_s1_q06_1}{SRG (Enthusiast) - Test 07 (27 Jun 2021)-img/p1_phy_s1_q06_1.png}{107.1}{0}
\qfigdef{p1_phy_s2_q07_1}{SRG (Enthusiast) - Test 07 (27 Jun 2021)-img/p1_phy_s2_q07_1.png}{206.0}{0}
\qfigdef{p1_phy_s2_q08_1}{SRG (Enthusiast) - Test 07 (27 Jun 2021)-img/p1_phy_s2_q08_1.png}{108.6}{0}
\qfigdef{p1_phy_s2_q09_1}{SRG (Enthusiast) - Test 07 (27 Jun 2021)-img/p1_phy_s2_q09_1.png}{184.2}{0}
\qfigdef{p1_phy_s2_q10_1}{SRG (Enthusiast) - Test 07 (27 Jun 2021)-img/p1_phy_s2_q10_1.png}{143.4}{0}
\qfigdef{p1_phy_s2_q12_1}{SRG (Enthusiast) - Test 07 (27 Jun 2021)-img/p1_phy_s2_q12_1.png}{274.3}{1}
\qfigdef{p1_phy_s3_q14_1}{SRG (Enthusiast) - Test 07 (27 Jun 2021)-img/p1_phy_s3_q14_1.png}{89.3}{0}
\qfigdef{p1_phy_s3_q17_1}{SRG (Enthusiast) - Test 07 (27 Jun 2021)-img/p1_phy_s3_q17_1.png}{150.7}{0}
\qfigdef{p1_phy_s3_q18_1}{SRG (Enthusiast) - Test 07 (27 Jun 2021)-img/p1_phy_s3_q18_1.png}{160.4}{0}
\qfigdef{p1_chem_s1_q01_1}{SRG (Enthusiast) - Test 07 (27 Jun 2021)-img/p1_chem_s1_q01_1.png}{87.7}{0}
\qfigdef{p1_chem_s1_q01_2}{SRG (Enthusiast) - Test 07 (27 Jun 2021)-img/p1_chem_s1_q01_2.png}{93.6}{0}
\qfigdef{p1_chem_s1_q01_3}{SRG (Enthusiast) - Test 07 (27 Jun 2021)-img/p1_chem_s1_q01_3.png}{82.9}{0}
\qfigdef{p1_chem_s1_q01_4}{SRG (Enthusiast) - Test 07 (27 Jun 2021)-img/p1_chem_s1_q01_4.png}{88.8}{0}
\qfigdef{p1_chem_s1_q02_1}{SRG (Enthusiast) - Test 07 (27 Jun 2021)-img/p1_chem_s1_q02_1.png}{175.2}{0}
\qfigdef{p1_chem_s1_q03_1}{SRG (Enthusiast) - Test 07 (27 Jun 2021)-img/p1_chem_s1_q03_1.png}{86.0}{0}
\qfigdef{p1_chem_s1_q03_2}{SRG (Enthusiast) - Test 07 (27 Jun 2021)-img/p1_chem_s1_q03_2.png}{84.5}{0}
\qfigdef{p1_chem_s1_q05_1}{SRG (Enthusiast) - Test 07 (27 Jun 2021)-img/p1_chem_s1_q05_1.png}{40.5}{0}
\qfigdef{p1_chem_s1_q05_A}{SRG (Enthusiast) - Test 07 (27 Jun 2021)-img/p1_chem_s1_q05_A.png}{47.8}{0}
\qfigdef{p1_chem_s1_q05_3}{SRG (Enthusiast) - Test 07 (27 Jun 2021)-img/p1_chem_s1_q05_3.png}{90.1}{0}
\qfigdef{p1_chem_s1_q05_C}{SRG (Enthusiast) - Test 07 (27 Jun 2021)-img/p1_chem_s1_q05_C.png}{63.5}{0}
\qfigdef{p1_chem_s1_q05_D}{SRG (Enthusiast) - Test 07 (27 Jun 2021)-img/p1_chem_s1_q05_D.png}{120.5}{0}
\qfigdef{p1_chem_s1_q06_B}{SRG (Enthusiast) - Test 07 (27 Jun 2021)-img/p1_chem_s1_q06_B.png}{59.1}{0}
\qfigdef{p1_chem_s1_q06_C}{SRG (Enthusiast) - Test 07 (27 Jun 2021)-img/p1_chem_s1_q06_C.png}{105.4}{0}
\qfigdef{p1_chem_s1_q06_D}{SRG (Enthusiast) - Test 07 (27 Jun 2021)-img/p1_chem_s1_q06_D.png}{91.5}{0}
\qfigdef{p1_chem_s2_q08_A}{SRG (Enthusiast) - Test 07 (27 Jun 2021)-img/p1_chem_s2_q08_A.png}{223.0}{0}
\qfigdef{p1_chem_s2_q08_B}{SRG (Enthusiast) - Test 07 (27 Jun 2021)-img/p1_chem_s2_q08_B.png}{194.4}{0}
\qfigdef{p1_chem_s2_q08_C}{SRG (Enthusiast) - Test 07 (27 Jun 2021)-img/p1_chem_s2_q08_C.png}{298.7}{1}
\qfigdef{p1_chem_s2_q08_D}{SRG (Enthusiast) - Test 07 (27 Jun 2021)-img/p1_chem_s2_q08_D.png}{225.6}{0}
\qfigdef{p1_chem_s2_q09_A}{SRG (Enthusiast) - Test 07 (27 Jun 2021)-img/p1_chem_s2_q09_A.png}{283.5}{1}
\qfigdef{p1_chem_s2_q09_B}{SRG (Enthusiast) - Test 07 (27 Jun 2021)-img/p1_chem_s2_q09_B.png}{205.3}{0}
\qfigdef{p1_chem_s2_q09_D}{SRG (Enthusiast) - Test 07 (27 Jun 2021)-img/p1_chem_s2_q09_D.png}{114.4}{0}
\qfigdef{p1_chem_s2_q11_1}{SRG (Enthusiast) - Test 07 (27 Jun 2021)-img/p1_chem_s2_q11_1.png}{178.3}{0}
\qfigdef{p1_chem_s2_q11_2}{SRG (Enthusiast) - Test 07 (27 Jun 2021)-img/p1_chem_s2_q11_2.png}{79.0}{0}
\qfigdef{p1_chem_s2_q12_A}{SRG (Enthusiast) - Test 07 (27 Jun 2021)-img/p1_chem_s2_q12_A.png}{174.7}{0}
\qfigdef{p1_chem_s2_q12_1}{SRG (Enthusiast) - Test 07 (27 Jun 2021)-img/p1_chem_s2_q12_1.png}{213.8}{0}
\qfigdef{p1_chem_s2_q12_C}{SRG (Enthusiast) - Test 07 (27 Jun 2021)-img/p1_chem_s2_q12_C.png}{166.0}{0}
\qfigdef{p1_chem_s2_q12_D}{SRG (Enthusiast) - Test 07 (27 Jun 2021)-img/p1_chem_s2_q12_D.png}{242.6}{0}
\qfigdef{p1_chem_s3_q13_1}{SRG (Enthusiast) - Test 07 (27 Jun 2021)-img/p1_chem_s3_q13_1.png}{144.9}{0}
\qfigdef{p1_chem_s3_q17_1}{SRG (Enthusiast) - Test 07 (27 Jun 2021)-img/p1_chem_s3_q17_1.png}{130.5}{0}
\qfigdef{p1_chem_s3_q18_1}{SRG (Enthusiast) - Test 07 (27 Jun 2021)-img/p1_chem_s3_q18_1.png}{298.9}{1}
\raggedcolumns\makeatletter\@beginparpenalty=10000\makeatother

\begin{document}
\begin{center}{\Large\bfseries SRG (Enthusiast) --- Test 07 (27 Jun 2021)}\end{center}
\begin{multicols}{2}
\partsec{1}{PHYSICS}
\sect{SECTION-I (i)}
\qn{1} A hemispherical non-conducting smooth bowl is kept fixed as shown in figure. A charge $Q$ is fixed at centre of bowl. A small block of mass $m$ and carrying charge $Q$ is released from the top of the bowl. It slides down and loses contact after covering angle $\theta$ as shown. If $\dfrac{KQ^2}{R^2}=\dfrac{mg}{2}$, then the value of $\cos\theta$ is :-
\qfig{p1_phy_s1_q01_1}
\begin{qoptions}
\opt{A}{$\dfrac23$}
\opt{B}{$\dfrac34$}
\opt{C}{$\dfrac45$}
\opt{D}{$\dfrac56$}
\end{qoptions}
\qn{2} From a charged, non-conducting, fixed thin disc of radius $R$ having total charge density $\sigma$, a disc of radius $\dfrac R2$ is cut and removed as shown in figure. The work done by external agent to slowly move a point charge $q$ from point $C$ to $D$ is :- [$C$ was the centre of original disc]
\qfig{p1_phy_s1_q02_1}
\begin{qoptions}
\opt{A}{$\dfrac{q\sigma R}{2\varepsilon_0}\left[\dfrac1\pi-\dfrac12\right]$}
\opt{B}{$\dfrac{q\sigma R}{\varepsilon_0}\left[\dfrac1\pi-\dfrac12\right]$}
\opt{C}{$\dfrac{q\sigma R}{\pi\varepsilon_0}\left[\dfrac1\pi-\dfrac12\right]$}
\opt{D}{$\dfrac{q\sigma R}{\varepsilon_0}\left[\dfrac12-\dfrac1\pi\right]$}
\end{qoptions}
\qn{3} A small section of area $\Delta A$ is removed from a spherical shell with surface mass density $\sigma$ and radius $R$ as shown in the figure. Find the magnitude of gravitational field intensity at point $P$ due to remaining mass :
\qfig{p1_phy_s1_q03_1}
\begin{qoptions}
\opt{A}{$\dfrac{4\sigma\Delta AG}{9R^2}$}
\opt{B}{$\dfrac{5\sigma\Delta AG}{9R^2}$}
\opt{C}{$\dfrac{\sigma\Delta AG}{9R^2}$}
\opt{D}{$\dfrac{7\sigma\Delta AG}{9R^2}$}
\end{qoptions}
\qn{4} The intensity of an electric field depends only on the co-ordinates $x$, $y$ and $z$ as follows:
{\small\[\vec E=a\,\dfrac{\left(x\hat i+y\hat j+z\hat k\right)}{\left(x^2+y^2+z^2\right)^{3/2}}\text{ unit.}\]}
The electrostatic energy stored between two imaginary concentric spherical shells of radii $R$ and $2R$ with centre at origin is :
\begin{qoptions}
\opt{A}{$\dfrac{4\pi\varepsilon_0a^2}{R}$}
\opt{B}{$\dfrac{2\pi\varepsilon_0a^2}{R}$}
\opt{C}{$\dfrac{\pi\varepsilon_0a^2}{R}$}
\opt{D}{$\dfrac{\pi\varepsilon_0a^2}{2R}$}
\end{qoptions}
\qn{5} A particle is projected from the surface of the earth (of radius $R_e$ and mass $M_e$) with speed equal to $\sqrt{\dfrac{GM_e}{R_e}}$, at certain angle from local horizontal as shown in the figure such that the angle subtended by arc between launching \& landing site at earth's centre is $\beta=60^\circ$. Calculate the maximum separation of the particle from the centre of the earth. Consider earth to be uniformly dense and air resistance to be absent.
\qfig{p1_phy_s1_q05_2}
\begin{qoptions}
\opt{A}{$R_e\left(1+\dfrac{\sqrt3}{2}\right)$}
\opt{B}{$R_e\left(\sqrt2+1\right)$}
\opt{C}{$R_e\dfrac{\left(\sqrt2+1\right)}{\sqrt2}$}
\opt{D}{$R_e\dfrac{\left(\sqrt2+1\right)}{\left(\sqrt2-1\right)}$}
\end{qoptions}
\qn{6} A light elastic cord of force constant $k$ and relaxed length $\ell$ is suspended from the ceiling. A bead of mass $m$ threaded on the cord is held motionless close to the ceiling and released. If the bead slides down the cord with a constant speed $v$ relative to the ceiling, find expression for thermal power dissipated during sliding of the bead.
\qfig{p1_phy_s1_q06_1}
\begin{qoptions}
\opt{A}{$mgv\left[\dfrac{mg+2k\ell}{2(mg+k\ell)}\right]$}
\opt{B}{$mgv\left[\dfrac{2mg+k\ell}{2(mg+k\ell)}\right]$}
\opt{C}{$mgv\left[\dfrac{mg+k\ell}{2(mg+k\ell)}\right]$}
\opt{D}{$mgv\left[\dfrac{mg+2k\ell}{mg+k\ell}\right]$}
\end{qoptions}
\sect{SECTION-I (ii)}
\qn{7} Consider a solid conducting sphere of radius $2R$. It has a spherical cavity of radius $R$ concentric with the solid sphere. The sphere is grounded. A point charge $q$ is fixed at point $P$ at distance of $5R$. Choose the correct statement(s) :
\qfig{p1_phy_s2_q07_1}
\begin{qoptions}
\opt{A}{The magnitude of charge on outer surface of solid sphere is $\left(\dfrac{2q}{5}\right)$C}
\opt{B}{The electrostatic energy density at point $P_1$ is zero.}
\opt{C}{The magnitude of net electric force experienced by $q$ due to charge on the sphere is $\left[\dfrac{2Kq}{125R^2}\right]$N/C}
\opt{D}{If the connection of sphere to ground is now removed then the potential energy of the system will increase}
\end{qoptions}
\qn{8} As shown in figure, a point charge $q_1$ is placed inside a neutral conducting shell at a distance $\dfrac R2$ from the centre of shell. Another point charge $q_2$ is placed at a distance $\dfrac{\sqrt7}{2}R$ from the centre of shell. Choose the correct option(s).
\qfig{p1_phy_s2_q08_1}
\begin{qoptions}
\opt{A}{The magnitude of force on $q_1$ due to induced charges at the inner surface of the conducting shell is zero.}
\opt{B}{The magnitude of force on $q_1$ due to induced charges at the outer surface of the conducting shell is $\dfrac{q_1q_2}{8\pi\epsilon_0R^2}$}
\opt{C}{The magnitude of force on $q_2$ due to induced charges at the inner surface of the shell is zero.}
\opt{D}{The magnitude of force on $q_2$ due to induced charges at the inner surface of the shell is $\dfrac{q_1q_2}{8\pi\epsilon_0R^2}$}
\end{qoptions}
\qn{9} Consider a system of two conducting spherical shells having common centre at $O$. The charge given to inner shell and outer shell is $Q_1$ and $Q_2$ respectively. A point charge $Q_4$ kept at distance $a$ from centre $O$. A point charge $Q_3$ is kept at distance $b$ from centre $O$. $\left[K=\dfrac1{4\pi\epsilon_0}\right]$. Choose the correct statement(s)
\qfig{p1_phy_s2_q09_1}
\begin{qoptions}
\opt{A}{The potential at centre $O$ is\\ {\small$\left[\dfrac{KQ_1}{R_1}+\dfrac{KQ_2}{R_2}+\dfrac{KQ_3}{b}+\dfrac{KQ_4}{a}\right]$}}
\opt{B}{The magnitude of electric field intensity at centre $O$ is $\left(\dfrac{KQ_4}{a^2}\right)$(N/C)}
\opt{C}{The magnitude of electric field intensity at centre $O$ due to charge on outer surface of outer shell is $\left[\dfrac{KQ_3}{b^2}\right]$}
\opt{D}{The magnitude of electric field intensity at point $P_2$ is $\dfrac{K(Q_1+Q_4)}{c^2}$}
\end{qoptions}
\columnbreak\qn{10} The charge per unit length of the four quadrants of the ring of radius $R$ is $2\lambda,-2\lambda,\lambda$ and $-\lambda$ respectively as shown in the figure
\qfig{p1_phy_s2_q10_1}
\begin{qoptions}
\opt{A}{The electric field at the centre is $\dfrac{\lambda}{2\pi\epsilon_0R}\left(-\hat i\right)$}
\opt{B}{The electric field at the centre is $\dfrac{\lambda}{2\pi\epsilon_0R}\left(-\hat j\right)$}
\opt{C}{The electric dipole moment of given charge distribution is $4\lambda R^2\left(\hat i\right)$}
\opt{D}{Electric potential at the centre of ring is zero}
\end{qoptions}
\qn{11} The electric potential $V$ in volts in a region of space is given by {\small\[V=ax^2+ay^2+2az^2\]} where $a$ is a constant of proper dimensions. Choose the correct statement(s)
\begin{qoptions}
\opt{A}{The work done by field when a $2\mu$C test charge moves from point $(0,0,0.1\text{ m})$ to origin is $5\times10^{-5}$ J. (The numeric value of $a$ is $1.25\times10^{3}$)}
\opt{B}{The work done by field when a $2\mu$C test charge moves from point $(0,0,0.1\text{ m})$ to origin is $5\times10^{-5}$ J. (The value of $a$ is $2.5\times10^{3}$)}
\opt{C}{The radius of the circle of equipotential line in $xy$-plane corresponding to $V=6250$V and $z=\sqrt2$m is 1 m. (The numeric value of $a$ is $1.25\times10^{3}$)}
\opt{D}{The radius of the circle of equipotential line in $xy$-plane corresponding to $V=6250$V and $z=\sqrt2$m is 0.1m. (The numeric value of $a$ is $1.25\times10^{3}$)}
\end{qoptions}
\qfullbegin
\qn{12} Figure shows a very long smooth non conducting rod of radius $r$, charged with uniform charge density $\lambda$, fixed horizontally.
\qfig{p1_phy_s2_q12_1}
\qfullsplit
A neutral and smooth ring $B$ of mass $M$ can slide freely on the rod which happens to just fit in it. $A$ is a nonconducting particle having mass $m$ and charge $q$, attached to the ring by means of a non-conducting and inextensible string of length $R$. If $A$ is released from the position shown in figure then :
\begin{qoptions}
\opt{A}{The loss in electrostatic potential energy when string becomes vertical is $\dfrac{q\lambda}{2\pi\varepsilon_0}\ell n\left[1+\dfrac Rr\right]$}
\opt{B}{The loss in electrostatic potential energy when string becomes vertical is $\dfrac{q\lambda}{4\pi\varepsilon_0}\ell n\left[1+\dfrac Rr\right]$}
\opt{C}{If $M=2m$ and $R=2r$, then speed of particle when string becomes vertical is $\left[\dfrac43\left(2gR+\dfrac{q\lambda}{\pi\varepsilon_0m}\ell n3\right)\right]^{\frac12}$}
\opt{D}{If $M=2m$ and $R=2r$, then speed of particle when string becomes vertical is $\left[\dfrac23\left(2gR+\dfrac{q\lambda}{\pi\varepsilon_0m}\ell n3\right)\right]^{\frac12}$}
\end{qoptions}
\qfullend
\sect{SECTION-II}
\qn{13} Three uniformly positively charged wires with linear charge density $\lambda$ are placed along $x$, $y$ and $z$ axis respectively. Find the flux of electric field through Gaussian surface given by $x^2+y^2+z^2=1$; $x>0$; $y>0$; $z>0$. If your answer is $\dfrac{n\lambda}{12\epsilon_0}$, fill the value of $n$.
\qn{14} A uniform cloud of dust particles approaches a planet from a great distance with a velocity $v_0$ relative to the planet. Radius of the planet is $R$ and escape velocity from the surface of the planet is $v_e$. Length of the cloud along its direction of motion is $\ell$, cross section of the cloud perpendicular to its direction of motion is very large and its density is $\rho$. How much mass of the dust particles will the planet collect during its passage through the cloud. Mass of the cloud is negligible as compared to the mass of the planet.
\qfig{p1_phy_s3_q14_1}
If your answer is $\dfrac{\pi\rho\ell R^{A}}{B}\left[1+\left(\dfrac{v_e}{v_0}\right)^{\frac C2}\right]$. Find $A+B+C$.
\qn{15} A heavy load is suspended on a light spring. The spring is slowly pulled down at the midpoint (a certain work $A$ is done thereby) and then released. The maximum kinetic energy $W_k$ of the load in the subsequent motion is $\dfrac An$. Find the value of $n$.
\qn{16} A solid non-conducting hemisphere of radius $R$ has a uniformly distributed positive charge of density $\rho$ per unit volume. A negatively charge particle having charge $q$ is transferred from centre of its base to infinity. Dielectric constant of material of hemisphere is unity. The external work performed in the process is found to be $\dfrac{q\rho R^2}{K\epsilon_0}$. The value of $K$ is
\qn{17} Two identical small balls $A$ and $B$ each of mass $m$ connected by a light inextensible cord of length $\ell$ are placed on a frictionless horizontal floor. With what minimum velocity $u$ must the ball $B$ be projected vertically upwards so that the ball $A$ leaves the floor ? Acceleration of free fall is $g$. If the answer is $\sqrt{ng\ell}$. Find $n$.
\qfig{p1_phy_s3_q17_1}
\qn{18} The system in figure is released from rest from the position shown. After blocks have moved distance $H/3$, collar $B$ is removed and block $A$ \& $C$ continue to move. The speed of $C$ just before it strikes the ground is $\dfrac{\sqrt{n\,gH}}{3}$. Find $n$
\qfig{p1_phy_s3_q18_1}
\partsec{2}{CHEMISTRY}
\sect{SECTION-I (i)}
\qn{1} Match list I with list II and select the \textbf{CORRECT} answer from the given below the lists

List I:

\noindent\begin{tabular}{@{}l@{\,}l@{\;\;}l@{\,}l@{}}
P. & \qfigi[0.95]{p1_chem_s1_q01_1} & Q. & \qfigi[0.95]{p1_chem_s1_q01_2}\\[4pt]
R. & \qfigi[0.95]{p1_chem_s1_q01_3} & S. & \qfigi[0.95]{p1_chem_s1_q01_4}
\end{tabular}

List II:

\noindent\begin{tabular}{@{}l@{\;}l@{}}
i. & $(2R,3E)-3-\text{pentene}-2-\text{ol}$\\
ii. & $(2R,3Z)-3-\text{pentene}-2-\text{ol}$\\
iii. & $(2S,3E)-3-\text{pentene}-2-\text{ol}$\\
iv. & $(2S,3Z)-3-\text{pentene}-2-\text{ol}$
\end{tabular}

\noindent\begin{tabular}{@{}l@{\quad}l@{\quad}l@{\quad}l@{\quad}l@{}}
 & P & Q & R & S\\
(A) & i & iii & iv & ii\\
(B) & i & iii & ii & iv\\
(C) & iv & iii & i & ii\\
(D) & iv & iii & ii & i
\end{tabular}
\qn{2} Select the \textbf{INCORRECT} statement for the compound
\qfig{p1_chem_s1_q02_1}
\begin{qoptions}
\opt{A}{Double bond has $Z$ configuration}
\opt{B}{Hydrogen deficiency index is $4$}
\opt{C}{Contains three stereogenic sites}
\opt{D}{No. of optical active isomers possible $=4$}
\end{qoptions}
\qn{3} Which one of the following statement is \textbf{CORRECT} about

\noindent\begin{tabular}{@{}l@{\,}l@{\;\;}c@{\;\;}l@{\,}l@{}}
 & \qfigi{p1_chem_s1_q03_1} & and & & \qfigi{p1_chem_s1_q03_2}\\
 & I & & & II
\end{tabular}
\begin{qoptions}
\opt{A}{I is optically inactive and II is optically active}
\opt{B}{$-$OH group occupies axial position is I, but equatorial position in II, in their most stable conformers}
\opt{C}{Ease of ionization of alcohol in acidic medium is I $>$ II}
\opt{D}{I is optically active and II is optically inactive}
\end{qoptions}
\qn{4} Poor shielding of nuclear charge by $d$ or $f$-orbital electrons is responsible for which of the following facts ?
\begin{qoptions}
\opt{A}{Atomic radius of Nb (4d-series) is greater than that of Ta (5d-series)}
\opt{B}{The 1$^{\text{st}}$ ionisation enthalpy of copper is less than that of zinc}
\opt{C}{The value of electron gain enthalpy is more negative for sulphur than oxygen.}
\opt{D}{The 1$^{\text{st}}$ ionisation energy for gold is greater than that of silver.}
\end{qoptions}
\qn{5} \qfigi{p1_chem_s1_q05_1} $\xrightarrow[\text{(monochlorination)}]{\text{Cl}_2/h\upsilon}$ major product $\xrightarrow{\text{Mg/dry ether}}$ P $\xrightarrow{\text{D}_2\text{O}}$ Q.

The product Q is
\begin{qoptions}
\opt{A}{\qfigi{p1_chem_s1_q05_A}}
\opt{B}{\qfigi{p1_chem_s1_q05_3}}
\opt{C}{\qfigi{p1_chem_s1_q05_C}}
\opt{D}{\qfigi{p1_chem_s1_q05_D}}
\end{qoptions}
\qn{6} {\small\[\mathrm{H_3C-\underset{\substack{|\\ \mathrm{CH_3}}}{CH}-CH_2-\underset{\substack{|\\ \mathrm{CH_3}}}{CH}-CH_2-CH_2-CH_3}\]\[\xrightarrow[\Delta]{\mathrm{Cr_2O_3/Al_2O_3}}\mathrm X\]}

The product $X$ is.
\begin{qoptions}
\opt{A}{$\mathrm{H_3C-\underset{\substack{|\\ \mathrm{CH_3}}}{\overset{\substack{\mathrm{CH_3}\\|}}{C}}-CH_2-\underset{\substack{|\\ \mathrm{CH_3}}}{\overset{\substack{\mathrm{CH_3}\\|}}{C}}-CH_3}$}
\opt{B}{\qfigi{p1_chem_s1_q06_B}}
\opt{C}{\qfigi{p1_chem_s1_q06_C}}
\opt{D}{\qfigi{p1_chem_s1_q06_D}}
\end{qoptions}
\sect{SECTION-I (ii)}
\qn{7} Which of the following is/are \textbf{CORRECT} statements about the following compound
{\small\[\mathrm{H_3C-\underset{\substack{|\\H}}{\overset{\substack{OH\\|}}{C}}-CH=CH-\underset{\substack{|\\H}}{\overset{\substack{OH\\|}}{C}}-CH_3}\]}
\begin{qoptions}
\opt{A}{Total no. of pair(s) of enantiomers $=2$}
\opt{B}{No. of optically inactive isomers $=2$}
\opt{C}{No. of optical isomers $=6$}
\opt{D}{No. of possible enantiomers, if double bond has trans configuration $=2$}
\end{qoptions}
\qfullbegin
\qn{8} Which of the following is/are \textbf{CORRECT} statement(s)
\begin{qoptions}
\opt{A}{\qfigi{p1_chem_s2_q08_A}}
\opt{B}{\qfigi{p1_chem_s2_q08_B}}
\opt{C}{\qfigi{p1_chem_s2_q08_C}}
\opt{D}{\qfigi{p1_chem_s2_q08_D}}
\end{qoptions}
\qn{9} Which of the following statement(s) is/are \textbf{CORRECT}.
\begin{qoptions}
\opt{A}{\qfigi{p1_chem_s2_q09_A}}
\opt{B}{\qfigi{p1_chem_s2_q09_B}\\ are asymmetric molecules and optically active}
\opt{C}{Specific rotation of $^{12}\mathrm C(\mathrm H)(\mathrm D)(\mathrm{CH_3})(\mathrm T)$ is greater than $^{14}\mathrm C(\mathrm H)(\mathrm D)(\mathrm{CH_3})(\mathrm T)$}
\opt{D}{\qfigi{p1_chem_s2_q09_D}}
\end{qoptions}
\qfullend
\qn{10} Which of the following statement(s) is/are \textbf{INCORRECT}.
\begin{qoptions}
\opt{A}{Acidified $\mathrm{H_2O_2}$ solution forms stable blue colour compound with $\mathrm{K_2Cr_2O_7}$ aqueous solution}
\opt{B}{Acidified $\mathrm{MnO_2}$ liberates $\mathrm{O_2}$ gas with $\mathrm{H_2O_2}$ solution}
\opt{C}{Alkaline $\mathrm{H_2O_2}$ solution turns $\mathrm{K_4\left[Fe(CN)_6\right]}$ solution into $\mathrm{K_3\left[Fe(CN)_6\right]}$}
\opt{D}{$\mathrm{PbO_2}$ is a stronger oxidizing agent than $\mathrm{H_2O_2}$ due to presence of weak peroxy link in $\mathrm{PbO_2}$.}
\end{qoptions}
\qn{11}
\qfig{p1_chem_s2_q11_1}
Which of the following is/are \textbf{CORRECT} statements
\begin{qoptions}
\opt{A}{Q is \qfigi{p1_chem_s2_q11_2}}
\opt{B}{stereomers of product R involves internal as well as external compensation of optical rotation}
\opt{C}{Single enantiomerically pure form of Q is formed}
\opt{D}{Enantiomers as well as diastereomers are formed as R}
\end{qoptions}
\qn{12} Which of the following is/are \textbf{CORRECT} ?
\begin{qoptions}
\opt{A}{\qfigi{p1_chem_s2_q12_A}}
\opt{B}{\qfigi{p1_chem_s2_q12_1}}
\opt{C}{\qfigi{p1_chem_s2_q12_C}}
\opt{D}{\qfigi{p1_chem_s2_q12_D}}
\end{qoptions}
\columnbreak\sect{SECTION-II}
\qn{13} In the following compound
\qfig{p1_chem_s3_q13_1}
If

\hspace*{1em}no. of optical active isomers $=X$

\hspace*{1em}no. of meso isomers \hspace{2.4em}$=Y$

\hspace*{1em}no. of isomers having pseudo chiral centers $=Z$

then, the value of $\left(\dfrac{X}{Y+Z}\right)$ is:
\qn{14} In the following statements about $\mathrm{H_2O_2}$
\begin{enumerate}[leftmargin=1.8em,itemsep=0pt,topsep=2pt,label=\arabic*.]
\item $\mathrm{H_2O_2}$ is diamagnetic in nature
\item Boiling point: $\mathrm{H_2O_2>H_2O}$ due to presence of strong H-bonds in $\mathrm{H_2O_2}$ than $\mathrm{H_2O}$
\item $\mathrm{Fe^{+2}}$ ion destabilises $\mathrm{H_2O_2}$ solution
\item Ethylene glycol stabilizes $\mathrm{H_2O_2}$ solution
\item Usually $\mathrm{H_2O_2}$ samples are stored in glass bottles
\item $\mathrm{H_2O_2}$ sample is used to clean old oil paintings
\item Sodium peroxo borate is a stable alternate source for $\mathrm{H_2O_2}$, which generates $\mathrm{H_2O_2}$ during cyclic anionic hydrolysis
\item $\mathrm{H_2O_2}$ turns acidified KI and starch paper into blue.
\item $\mathrm{H_2O_2}$ liberates $\mathrm{N_2}$ gas with $\mathrm{N_2H_4}$ and this reaction is exothermic.
\end{enumerate}
If:

no. of correct statements $=X$

no. of incorrect statements $=Y$

then the value of $(Y-X)$ is
\qn{15} In the following case

$\mathrm{O_3,Ag_2O,KMnO_4}$ (neutral medium),

$\mathrm{PbS,Cl_2,Ce^{+3}}$ (acidic medium),\\ $\mathrm{Au^{+3},NaOBr,NaOH,PbO_2,BaSO_3}$

If

No. of cases, in which $\mathrm{H_2O_2}$ as oxidising agent $=X$

No. of cases, in which $\mathrm{H_2O_2}$ as reducing agent $=Y$

then the value of $\left(\dfrac{X+Y}{X-Y}\right)$ is
\qn{16} In the following compounds

$\mathrm{BiBr_5,TlI_3,PbCl_4,SnCl_4,PbO_2,BiF_5,PbI_4,BiI_3,PbBr_2}$

If :

No. of compounds exist at room temperature $=X$

No. of compounds donot exist at room temperature $=Y$

then the product of $(XY)$ is
\qn{17} In the following mono chlorination reaction.
\qfig{p1_chem_s3_q17_1}
(enantiomerically pure)

If

No. of optically active products formed $=X$

No. of meso products formed $=Y$

possible no. of products $=Z$

then, the value of $\left(\dfrac{Y+Z}{X}\right)$ is
\qfullbegin
\qn{18} In the following reactions
\qfig{p1_chem_s3_q18_1}
\qfullsplit
no. of possible stereo isomers of $P$ formed in this reaction $=Y$ then, the value of $\left(\dfrac{Y+X}{Y-X}\right)$ is
\qfullend
\partsec{3}{MATHEMATICS}
\sect{SECTION-I (i)}
\qn{1} Let $f:(-\infty,-1]\to\left(\dfrac\pi2,\pi\right]$ be defined as $f(x)=\sec^{-1}(-x^2+x+a)$. If $f(x)$ is surjective then the range of $a$ is :
\begin{qoptions}
\opt{A}{$\{1\}$}
\opt{B}{$\left\{\dfrac{-5}{4}\right\}$}
\opt{C}{$\left(-\infty,\dfrac{-5}{4}\right]$}
\opt{D}{$(-\infty,1]$}
\end{qoptions}
\qn{2} The sum $\dfrac{19}{1\cdot2\cdot3\cdot4}+\dfrac{28}{2\cdot3\cdot4\cdot8}+\dfrac{39}{3\cdot4\cdot5\cdot16}+\dfrac{52}{4\cdot5\cdot6\cdot32}+\ldots+$ upto infinite terms is equal to
\begin{qoptions}
\opt{A}{$\dfrac12$}
\opt{B}{$1$}
\opt{C}{$2$}
\opt{D}{$4$}
\end{qoptions}
\qn{3} $\displaystyle\lim_{x\to0}\left(\dfrac{2}{x^3}\left(\sin^{-1}x-\tan^{-1}x\right)\right)^{\frac{2}{x^2}}=$
\begin{qoptions}
\opt{A}{$e$}
\opt{B}{$\sqrt e$}
\opt{C}{$\dfrac1e$}
\opt{D}{$\dfrac1{\sqrt e}$}
\end{qoptions}
\qn{4} If {\small\[\begin{aligned}\lim_{x\to0}\Bigg(&\left[\dfrac{\sin^{-1}x}{x}\right]+\left[\dfrac{2^2\sin^{-1}2x}{x}\right]\\&+\left[\dfrac{3^2\sin^{-1}3x}{x}\right]+\ldots+\left[\dfrac{n^2\sin^{-1}nx}{x}\right]\Bigg)=100\end{aligned}\]} then the value of $n$ is

[Note: $[K]$ denotes the greatest integer less than or equal to $k$]
\begin{qoptions}
\opt{A}{$2$}
\opt{B}{$3$}
\opt{C}{$4$}
\opt{D}{$5$}
\end{qoptions}
\qn{5} If $\displaystyle\lim_{x\to0}\dfrac{x+\sin x-x\cos x-\tan x}{x^n}$ exists and is non-zero finite value, then the value of $n$ is
\begin{qoptions}
\opt{A}{$3$}
\opt{B}{$4$}
\opt{C}{$5$}
\opt{D}{$6$}
\end{qoptions}
\qn{6} If $f(x)=\tan^{-1}x-\dfrac2\pi(\tan^{-1}x)^2+\dfrac4{\pi^2}(\tan^{-1}x)^3-\ldots\infty$, then the sum of integral values of $a$ for which equation $f^{\,2}(x)+(\sin^{-1}x)^2=a$, possess solution(s), is
\begin{qoptions}
\opt{A}{$6$}
\opt{B}{$7$}
\opt{C}{$9$}
\opt{D}{$10$}
\end{qoptions}
\sect{SECTION-I (ii)}
\qn{7} Let $\displaystyle\lim_{x\to\infty}(2^x+a^x+e^x)^{1/x}=L$, which of the following statement(s) is (are) correct?
\begin{qoptions}
\opt{A}{If $L=a(a>0)$, then the range of $a$ is $[e,\infty)$}
\opt{B}{If $L=2e(a>0)$, then the range of $a$ is $\{2e\}$}
\opt{C}{If $L=e(a>0)$, then the range of $a$ is $(0,e]$}
\opt{D}{If $L=e(a>0)$, then the range of $a$ is $[e,\infty)$}
\end{qoptions}
\qn{8} If range of the function {\small\[f(x)=\left(\cos^{-1}\dfrac x2\right)^2+\pi\sin^{-1}\dfrac x2-\left(\sin^{-1}\dfrac x2\right)^2+\dfrac{\pi^2}{12}(x^2+6x+8)\]} is $[a\pi^2,b\pi^2]$, then which of the following options are correct
\begin{qoptions}
\opt{A}{$a=\dfrac12$}
\opt{B}{$a=\dfrac14$}
\opt{C}{$a+b=\dfrac52$}
\opt{D}{$b-a=3$}
\end{qoptions}
\qn{9} Let $x_1,x_2,\ldots,x_{2018}$ be positive real numbers such that $x_1+x_2+\ldots+x_{2018}=1$. Possible value(s) of $k$ such that {\small\[k\sum_{i=1}^{2018}\dfrac{x_i^2}{1-x_1}\ge1\]} is (are):
\begin{qoptions}
\opt{A}{$2017$}
\opt{B}{$2018$}
\opt{C}{$2019$}
\opt{D}{$2020$}
\end{qoptions}
\qn{10} If $f(n)=\cot^{-1}(n+3)-2\cot^{-1}(n+1)+\cot^{-1}(n-1)$, $n\in N$, then $\displaystyle\sum_{n=1}^{\infty}f(n)$ is equal to;
\begin{qoptions}
\opt{A}{$\displaystyle\sum_{k=1}^{3}\tan^{-1}\left(\dfrac1k\right)$}
\opt{B}{$\displaystyle\lim_{x\to\pi}\dfrac{-\sin(x)}{\pi-x}$}
\opt{C}{$\displaystyle\lim_{x\to0^-}\cot^{-1}\left(\dfrac1x\right)$}
\opt{D}{$\displaystyle\lim_{x\to\pi}\dfrac{\pi(x-\pi)^2}{4(\cos x+1)}$}
\end{qoptions}
\qn{11} Let $f(x)=\dfrac{1-\cos\{x\}}{(x^4+ax^3+bx^2+cx)^2}$. If $l=\displaystyle\lim_{x\to1^+}f(x)$, $m=\displaystyle\lim_{x\to2^+}f(x)$ and $n=\displaystyle\lim_{x\to3^+}f(x)$ where $l$, $m$ and $n$ are non-zero finite then: (Note: $\{y\}$ denotes fractional part of $y$)
\begin{qoptions}
\opt{A}{$\displaystyle\lim_{x\to0^+}f(x)=\dfrac1{36}$}
\opt{B}{$a+b+c=-1$}
\opt{C}{$l+m+n=\dfrac{19}{72}$}
\opt{D}{$\displaystyle\lim_{x\to4^+}f(x)=0$}
\end{qoptions}
\qn{12} $\displaystyle\lim_{x\to\infty}2\left(\sqrt{25x^2+x}-5x\right)$ is equal to
\begin{qoptions}
\opt{A}{$\displaystyle\lim_{x\to0}\dfrac{2x-\log_e(1+x)^2}{5x^2}$}
\opt{B}{$\displaystyle\lim_{x\to0}\dfrac{e^{-x}-1+x}{x^2}$}
\opt{C}{$\displaystyle\lim_{x\to0}\dfrac{2(1-\cos x^2)}{5x^4}$}
\opt{D}{$\displaystyle\lim_{x\to0}\dfrac{\sin\frac x5}{x}$}
\end{qoptions}
\sect{SECTION-II}
\qn{13} If $\alpha$ and $\beta(\alpha<\beta)$ are the roots of the equation
{\small\[\lim_{t\to\infty}\cos^{-1}\left(\sin\left(\tan^{-1}\left(\dfrac{\sqrt{tx}}{\sqrt{tx^2-3tx+t-1}-x}\right)\right)\right)=\dfrac\pi6\]}
then find the value of $(8^{\alpha}+2^{\beta}-\alpha\beta)$
\qn{14} The sequence $\{a_n\}_{n=1}^{+\infty}$ is defined by $a_1=0$ and $a_{n+1}=a_n+4n+3$, $n\in N$. Find the value of
{\small\[\lim_{n\to+\infty}\dfrac{\sqrt{a_n}+\sqrt{a_{4n}}+\sqrt{a_{4^2n}}+\sqrt{a_{4^3n}}+\ldots+\sqrt{a_{4^{10}n}}}{\sqrt{a_n}+\sqrt{a_{2n}}+\sqrt{a_{2^2n}}+\sqrt{a_{2^3n}}+\ldots+\sqrt{a_{2^{10}n}}}\]}
\qn{15} Let {\small\[L=\lim_{n\to\infty}n\left(\sqrt[3]{n^3+3n^2+2n+1}+\sqrt{n^2-2n+3}-2n\right)\]} If $L$ can be expressed in the form of $\dfrac pq$ (where $p$ and $q$ are coprime) then find the value of $(p+q)$
\qn{16} If $\displaystyle\sum_{r=1}^{10}\tan^{-1}\left(\dfrac{3}{9r^2+3r-1}\right)=\cot^{-1}\left(\dfrac mn\right)$ (where $m$ and $n$ are coprime), then find $(2m+n)$
\qn{17} If $x_1$, $x_2$ and $x_3$ are the positive roots of the equation $x^3-6x^2+3px-2p=0$, $p\in R$. If the value of {\small\[\sin^{-1}\left(\dfrac1{x_1}+\dfrac1{x_2}\right)+\cos^{-1}\left(\dfrac1{x_2}+\dfrac1{x_3}\right)-\tan^{-1}\left(\dfrac1{x_3}+\dfrac1{x_1}\right)\]} is equal to $k$, then $\dfrac{12k}{\pi}$ is:
\qn{18} If {\small\[f(x)=\sin^{-1}\left(\dfrac{x^2(2x^2+1)-1}{1+x^2}\right)\]\[=\begin{cases}a\pi+b\cos^{-1}x,&0\le x\le1\\ p\pi+q\sin^{-1}x,&-1\le x<0\end{cases}\]} then find the absolute value of $(a+b+p+q)$.
\end{multicols}
\end{document}
